\section{Appendix: moment matching and proof details}

This appendix records the dense prior construction that supplies the superquadratic-sample lower bound invoked in \cref{obj:thm:all-estimator-lower}. The construction embeds many small, equal-mass treatment-effect contrasts into the observed-law class, then applies the classical moment-matching method for \(|t|\) \citep{CaiLow2011L1} through a Poissonized observation experiment of the type used in large-alphabet \(L_1\) lower bounds \citep{JiaoHanWeissman2018L1}. The role of the appendix is to identify the objects entering that reduction and state the quantitative lower-bound lemma used by the main text.

The dense construction begins with two equivalent descriptions of the same moment-matching experiment. The first packages the ingredients as a single experiment, while the second expands the degree, amplitude, priors, mixtures, risk, and separation that enter the lower-bound calculation.

\begin{definitionv}{synth_10}[Dense moment-matching construction]
For \(n,d\in\mathbb N\), \(0<\epsilon<1/2\), a domain certificate \(2\le d\) and \(0<a_{n,d}\le 1/2\), and a dense prior family \(K\mapsto(\nu_{0,K},\nu_{1,K})\) satisfying the stated moment-matching certificate at every positive even degree, define \(\operatorname{DenseMean}_{n,d,\epsilon}\) to be the dense moment-matching experiment consisting of the dense laws \(P_\theta^{\mathrm{dense}}\), the lower-bound degree \(K_d^{\mathrm{lb}}=2\lceil4L_d\rceil\), the intensity \(\lambda_{n,d}=2n/d\), the amplitude \(a_{n,d}=\sqrt{K_d^{\mathrm{lb}}/(64\lambda_{n,d})}\), the coordinate priors \(\nu_{0,K_d^{\mathrm{lb}}}\) and \(\nu_{1,K_d^{\mathrm{lb}}}\), their scaled product priors \(\Pi_{0,n,d}\) and \(\Pi_{1,n,d}\) on \([-a_{n,d},a_{n,d}]^d\), the induced Poissonized observation mixtures \(\mathbb M_{0,n,d}\) and \(\mathbb M_{1,n,d}\) with independent total sample size \(\operatorname{Pois}(2n)\), the Poissonized risk \(\mathfrak R^{\mathrm{Pois}}_{2n,d,\epsilon}\), and the prior target-mean separation \(\Delta_{n,d}^{\mathrm{dense}}=E_{\Pi_{1,n,d}}\Psi\{\operatorname{Obs}(P_\theta^{\mathrm{dense}})\}-E_{\Pi_{0,n,d}}\Psi\{\operatorname{Obs}(P_\theta^{\mathrm{dense}})\}\).
\end{definitionv}

\begin{definitionv}{def:dense-moment-matching-construction}[Dense moment matching]
For \(n,d\in\mathbb N\), \(\epsilon\in\mathbb R\), a dense domain certificate
\[
2\le d,\qquad
0<\epsilon,\qquad
\epsilon<\frac12,\qquad
0<a_{n,d}\le \frac12,
\]
and a dense prior family \(K\mapsto(\nu_{0,K},\nu_{1,K})\) with its moment-matching certificate at every positive even degree \(K\), the dense moment-matching construction is the object with the following components:
\begin{itemize}
\item \textbf{(Dense law.)} For each \(\theta\in[-1/2,1/2]^d\), \(P_\theta^{\mathrm{dense}}\) is the potential-outcome probability law induced by the dense full-data cell masses.
\item \textbf{(Approximation scale.)}
\[
K_d^{\mathrm{lb}}=2\lceil 4L_d\rceil,\qquad
E_K=\inf\left\{e\in\mathbb R:\exists p,\ \deg(p)\le K\ \text{and}\ 
\forall t\in[-1,1],\ \bigl||t|-p(t)\bigr|\le e\right\}.
\]
\item \textbf{(Poisson intensity and amplitude.)}
\[
\lambda_{n,d}=\frac{2n}{d},\qquad
a_{n,d}=\sqrt{\frac{K_d^{\mathrm{lb}}}{64\lambda_{n,d}}}.
\]
\item \textbf{(Coordinate priors.)} The two coordinate priors are
\[
\nu_{0,K_d^{\mathrm{lb}}},\qquad \nu_{1,K_d^{\mathrm{lb}}}.
\]
\item \textbf{(Dense product priors.)} The priors \(\Pi_{0,n,d}\) and \(\Pi_{1,n,d}\) are obtained by restricting the product measures \(\nu_{0,K_d^{\mathrm{lb}}}^{\otimes d}\) and \(\nu_{1,K_d^{\mathrm{lb}}}^{\otimes d}\) to \([-1,1]^d\), scaling each coordinate by \(a_{n,d}\), and mapping the result into the dense contrast class.
\item \textbf{(Observation mixtures.)} The mixtures \(\mathbb M_{0,n,d}\) and \(\mathbb M_{1,n,d}\) are the observation-kernel mixtures induced by \(\Pi_{0,n,d}\) and \(\Pi_{1,n,d}\), with independent Poissonized sample size \(\operatorname{Pois}(2n)\).
\item \textbf{(Risk and separation.)} The risk component is \(\mathfrak R^{\mathrm{Pois}}_{2n,d,\epsilon}\), and the prior target-mean separation component is
\[
\Delta_{n,d}^{\mathrm{dense}}
=
\operatorname{DenseMean}_{n,d,\epsilon}(\nu_{1,K_d^{\mathrm{lb}}})
-
\operatorname{DenseMean}_{n,d,\epsilon}(\nu_{0,K_d^{\mathrm{lb}}}).
\]
\end{itemize}
The construction also carries the supplied prior-pair certificate at degree \(K_d^{\mathrm{lb}}\) and the identities
\[
K=K_d^{\mathrm{lb}},\qquad
\Delta_{n,d}^{\mathrm{dense}}
=
a_{n,d}E_{K_d^{\mathrm{lb}}}.
\]
\end{definitionv}

\Cref{obj:synth_10,obj:def:dense-moment-matching-construction} specify the lower-bound scale. The even degree \(K_d^{\mathrm{lb}}\) is logarithmic in the alphabet size, the intensity \(\lambda_{n,d}\) is the expected count per cell under a \(\operatorname{Pois}(2n)\) sample, and the amplitude \(a_{n,d}\) balances the moment-matching tail against the target separation. The identity \(\Delta_{n,d}^{\mathrm{dense}}=a_{n,d}E_{K_d^{\mathrm{lb}}}\), together with the approximation behavior \(E_K\asymp K^{-1}\) from \citet{CaiLow2011L1}, is the source of the \(d/\{n\log(ed)\}\) squared-risk scale in the nonsaturated regime.

The observation experiment and likelihood calculation are local across cells. The following definitions isolate the Poissonized sampling kernel, the likelihood tail left after matching moments through degree \(K_d^{\mathrm{lb}}\), and the one-cell likelihood ratio used to control the total variation between the two mixtures.

\begin{definitionv}{synth_8}[Dense Poisson observation kernel]
For \(n,d\in\mathbb N\) with \(d\ge2\), \(\operatorname{Obs}_{n,d}\) is the Markov kernel from dense contrasts \(\theta\in[-1/2,1/2]^d\) to finite observed samples that maps \(\theta\) to the law of an independent \(\operatorname{Pois}(2n)\) number of observations from \(\operatorname{Obs}(P_\theta^{\mathrm{dense}})\). Thus, for any prior on \(\theta\), composing the prior with \(\operatorname{Obs}_{n,d}\) gives the corresponding Poissonized observed-sample mixture.
\end{definitionv}

\begin{definitionv}{synth_11}[Dense likelihood tail]
For the dense lower-bound degree \(K_d^{\mathrm{lb}}\), Poissonized cell intensity \(\lambda_{n,d}=2n/d\), and dense amplitude \(a_{n,d}\), define the dense likelihood tail by \[\operatorname{tail}_{n,d}^{\mathrm{dense}}:=\sum_{r=K_d^{\mathrm{lb}}+1}^{\infty}\frac{(\lambda_{n,d}a_{n,d}^2)^r}{r!}.\]
\end{definitionv}

\begin{definitionv}{synth_9}[One-cell dense likelihood ratio]
Fix a Poisson cell intensity \(\lambda\ge0\). Let \(R_+\) and \(R_-\) be the two sign counts in one dense cell, and let \(E_0\) denote expectation under the zero-contrast baseline law \(R_+,R_-\stackrel{\mathrm{ind}}{\sim}\operatorname{Pois}(\lambda/2)\). For \(\theta\in[-1,1]\), the one-cell likelihood ratio relative to this baseline is \[ L_\theta(R_+,R_-):=(1+\theta)^{R_+}(1-\theta)^{R_-}. \] In the dense lower bound, \(L_{\theta_x}\) applies this definition with the cell contrast \(\theta_x\) and \(\lambda=\lambda_{n,d}\).
\end{definitionv}

\Cref{obj:synth_8} gives the mixture laws to which Le Cam's fuzzy-hypothesis method applies \citep{LeCam1986,Tsybakov2009}. \Cref{obj:synth_11,obj:synth_9} provide the analytic bookkeeping for that comparison: moment matching cancels the low-degree likelihood terms, and the remaining tail controls the total variation between the two product mixtures.

The quantitative lower-bound statement combines the dense observed submodel, the moment-matched priors, concentration of the target under those priors, and the depoissonization comparison to fixed sample size.

\begin{lemmav}{lem:dense-moment-matching-lower}[Dense moment matching lower bound]*
Suppose the product experiment satisfies the i.i.d. sampling law in \cref{obj:ass:iid-sampling}.  There is an integer cutoff \(D_0\geq2\) such that, for every overlap constant \(\epsilon\in(0,1/2)\), there is a constant \(c_\epsilon>0\) with the following property.  For all integers \(d,n\) satisfying \(d\geq D_0\) and \(d^2<n\):
\begin{itemize}
\item \textbf{(Degree and amplitude.)} The lower-bound degree and amplitude obey
\[
8L_d\leq K_d^{\mathrm{lb}}<8L_d+2,
\qquad
a_{n,d}^2=\frac{K_d^{\mathrm{lb}}d}{128n},
\qquad
a_{n,d}\leq \frac12 .
\]
\item \textbf{(Dense observed submodel.)} For every \(\theta\in[-a_{n,d},a_{n,d}]^d\), the dense equal-mass contrast submodel \(P_\theta^{\mathrm{dense}}\) has observed margin in \(\mathcal D_{d,\epsilon}^{\mathrm{obs}}\) and satisfies, for every \(x\in[d]\),
\[
p_x=\frac1d,
\qquad
\pi_x=\frac12,
\qquad
\mu_{1x}=\frac{1+\theta_x}{2},
\qquad
\mu_{0x}=\frac{1-\theta_x}{2},
\]
with
\[
\Psi\{\operatorname{Obs}(P_\theta^{\mathrm{dense}})\}
=
\frac12+\frac{1}{2d}\sum_{x=1}^d|\theta_x|.
\]
\item \textbf{(Moment-matched priors and mixtures.)} The two dense product priors \(\Pi_{0,n,d}\) and \(\Pi_{1,n,d}\), with induced Poissonized mixtures \(\mathbb M_{0,n,d}\) and \(\mathbb M_{1,n,d}\), may be chosen so that their prior-mean separation satisfies
\[
\Delta_{n,d}^{\mathrm{dense}}
=
a_{n,d}E_{K_d^{\mathrm{lb}}},
\qquad
0<\Delta_{n,d}^{\mathrm{dense}}\leq1,
\]
and
\[
\operatorname{TV}(\mathbb M_{0,n,d},\mathbb M_{1,n,d})
\leq
d\{\operatorname{tail}_{n,d}^{\mathrm{dense}}\}^{1/2}
\leq
d\left\{\frac{(e/64)^{K_d^{\mathrm{lb}}+1}}{1-e/64}\right\}^{1/2}
\leq\frac1{16}.
\]
Under each of the two product priors, the target is concentrated around its prior mean:
\[
\int
\left(\Psi\{\operatorname{Obs}(P_\theta^{\mathrm{dense}})\}
-\mathbb E_{\Pi_{j,n,d}}\Psi\{\operatorname{Obs}(P_\theta^{\mathrm{dense}})\}\right)^2
\,d\Pi_{j,n,d}(\theta)
\leq
\frac{a_{n,d}^2}{4d},
\qquad j\in\{0,1\},
\]
and
\[
\Pi_{j,n,d}\left(
\left|\Psi\{\operatorname{Obs}(P_\theta^{\mathrm{dense}})\}
-\mathbb E_{\Pi_{j,n,d}}\Psi\{\operatorname{Obs}(P_\theta^{\mathrm{dense}})\}\right|
>
\frac{\Delta_{n,d}^{\mathrm{dense}}}{4}
\right)
\leq
\frac18,
\qquad j\in\{0,1\}.
\]
\item \textbf{(One-cell likelihood identity.)} For all contrasts \(\theta,\theta'\in\mathbb R\),
\[
E_0[L_\theta L_{\theta'}]
=
\exp(\lambda_{n,d}\theta\theta').
\]
\end{itemize}
Consequently,
\[
c_\epsilon(\Delta_{n,d}^{\mathrm{dense}})^2
\leq
\mathfrak R^{\mathrm{Pois}}_{2n,d,\epsilon},
\qquad
c_\epsilon\frac{d}{nK_d^{\mathrm{lb}}}
\leq
\mathfrak R^{\mathrm{Pois}}_{2n,d,\epsilon},
\]
and the fixed-sample and Poissonized risks satisfy
\[
\mathfrak R_{n,d,\epsilon}
\geq
\mathfrak R^{\mathrm{Pois}}_{2n,d,\epsilon}
-
\Pr\{\operatorname{Pois}(2n)<n\}.
\]
In particular,
\[
c_\epsilon\frac{d}{nL_d}
\leq
\mathfrak R_{n,d,\epsilon}.
\]
\end{lemmav}
% CAUSALSMITH-CITED-SCOPE-BEGIN lem:dense-moment-matching-lower
\verificationfootnotetext{\textbf{Formalization scope.} This result uses the published conclusion \citep{CaiLow2011L1} (Lemma 1 and Section 3.1). The cited source proof is not formalized here: Lean verifies its use and all remaining steps. If that cited conclusion is false, the portions of this result that depend on it are not certified.}
% CAUSALSMITH-CITED-SCOPE-END lem:dense-moment-matching-lower

\Cref{obj:lem:dense-moment-matching-lower} supplies the dense-regime ingredient for the main lower bound. In the submodel, each cell has mass \(1/d\), propensity \(1/2\), and a symmetric treatment contrast, so the optimal-regression value equals a constant plus the average absolute contrast. The two priors from \citet{CaiLow2011L1} separate this average absolute contrast by \(\Delta_{n,d}^{\mathrm{dense}}\), while the displayed likelihood-tail bound keeps the induced Poissonized mixtures close enough for the fuzzy-hypothesis argument \citep{LeCam1986,Tsybakov2009}. Since \(K_d^{\mathrm{lb}}\) is proportional to \(L_d\), the lemma yields the fixed-sample lower scale \(d/(nL_d)\) in the range \(d^2<n\), complementing the \(L_1\) embedding argument used in \cref{obj:thm:all-estimator-lower}.
